\scanpage{26}
des fractions de $A$. Comme c'est un module de
type fini sur $A$, il est dans la clôture
intégrale de $A$, O.K.

\uline{Question} Les conditions du cor. 2 sont elles
\uline{tjrs} vérifiées ? On est ramené évid. au cas
où les idéaux pr. associés à $0$ dans $A$
sont de \deleted{dim} corang $n$ : est-il vrai alors
que $H^n(\Omega)$ est dualisant ? \uuline{Non} déjà
dans le cas \deleted{d'une surface \unsure{non normale} ;} \inserted{\unsure{voir} \deleted{\unsure{l'ex.}} un anneau}
local \inserted{intègre} $A$ [de dim 2], tel que $A_p$ soit
régulier pour $p$ \deleted{\unsure{max}}\inserted{premier \unsure{minimal|minimum}} [$A$ non normal]. En
effet, \deleted{\unsure{on vérifie que}} on prouve alors que
\[
\operatorname{Hom}_A(\Omega,\Omega)=\bigcap_{\substack{p\ \text{premier}\\\text{de rang }1}}A_p
=\deleted{\text{\unsure{normalisé}}}\inserted{\text{clôture intégrale}}\text{ de }A .
\]
[$\Omega$ est en effet défini par des conditions \unsure{valuatives}
\mbox{-{}-{}-} ].

\uline{Proposition 10} Soit $x$ un élément de $A$ qui
n'y est pas diviseur de $0$, et tel que
$xH^{n-1}(A)=H^{n-1}(A)$. Alors $\Omega_A/x\Omega_A$ est
un module fondamental pour $\deleted{\unsure{\Omega_A|\Omega}}A/xA$
% archive sheet number: 25

